
\documentclass{article}
\usepackage{colortbl}
\usepackage{makecell}
\usepackage{multirow}
\usepackage{supertabular}

\begin{document}

\newcounter{utterance}

\centering \large Interaction Transcript for game `cladder', experiment `full\_v1.5\_default', episode 6543 with Qwen3.5{-}9B.
\vspace{24pt}

{ \footnotesize  \setcounter{utterance}{1}
\setlength{\tabcolsep}{0pt}
\begin{supertabular}{c@{$\;$}|p{.15\linewidth}@{}p{.15\linewidth}p{.15\linewidth}p{.15\linewidth}p{.15\linewidth}p{.15\linewidth}}
    \# & \multicolumn{2}{c}{Player} && \multicolumn{2}{c}{Game Master} \\
    \hline

    \theutterance \stepcounter{utterance}  
    & & & \multicolumn{4}{p{0.6\linewidth}}{
        \cellcolor[rgb]{0.9,0.9,0.9}{
            \makecell[{{p{\linewidth}}}]{
                \texttt{\tiny{[P1$\langle$GM]}}
                \texttt{You are an expert in causal inference. The following question is not a typical commonsense query, but rather a meticulously designed question created by a professor specializing in causal inference, intended to assess the students' mastery of the course content.} \\
\texttt{Imagine a self{-}contained, hypothetical world with only the following conditions, and without any unmentioned factors or causal relationships: Kwox has a direct effect on muvy and swoq. Swoq has a direct effect on kwoz. Muvy has a direct effect on kwoz. For those who are not kwox, the probability of kwoz is 34\%. For those who are kwox, the probability of kwoz is 55\%. Will kwox increase the chance of kwoz?} \\
\texttt{Start your answer with 'Yes' or 'No', followed by additional reasoning or evidence to support your explanation.} \\
            }
        }
    }
     \\ \\

    \theutterance \stepcounter{utterance}  
    & \multicolumn{4}{p{0.6\linewidth}}{
        \cellcolor[rgb]{0.561,0.737,0.561}{
            \makecell[{{p{\linewidth}}}]{
                \texttt{\tiny{[P1$\rangle$GM]}}
                \texttt{Yes} \\
\\ 
\texttt{**Reasoning:**} \\
\\ 
\texttt{To determine if \$Kwox\$ increases the chance of \$Kwoz\$, we must analyze the causal structure and the provided conditional probabilities within the constraints of this hypothetical world.} \\
\\ 
\texttt{1.  **Causal Structure**:} \\
\texttt{    *   \$Kwox \textbackslash{}rightarrow Muvy\$} \\
\texttt{    *   \$Kwox \textbackslash{}rightarrow Swoq\$} \\
\texttt{    *   \$Swoq \textbackslash{}rightarrow Kwoz\$} \\
\texttt{    *   \$Muvy \textbackslash{}rightarrow Kwoz\$} \\
\\ 
\texttt{    This structure indicates that \$Kwox\$ is a common cause (a confounder) for both \$Muvy\$ and \$Swoq\$, which in turn are direct causes of \$Kwoz\$. There are no other variables mentioned that influence \$Kwoz\$ or act as mediators between \$Kwox\$ and its effects.} \\
\\ 
\texttt{2.  **Causal Effect Definition**:} \\
\texttt{    The question asks if \$Kwox\$ *increases* the chance of \$Kwoz\$. In causal inference, this refers to the **Average Causal Effect (ACE)**, defined as the difference in the probability of the outcome (\$Kwoz\$) under intervention versus no intervention:} \\
\texttt{    \$\$ACE = P(Kwoz | do(Kwox)) {-} P(Kwoz | do(\textbackslash{}neg Kwox))\$\$} \\
\\ 
\texttt{    Since \$Kwox\$ is the root cause in this system and there are no unmentioned confounders or selection biases, the probability of \$Kwoz\$ given an intervention on \$Kwox\$ is equivalent to the observed conditional probability \$P(Kwoz | Kwox)\$. Similarly, the probability under no intervention is \$P(Kwoz | \textbackslash{}neg Kwox)\$.} \\
\\ 
\texttt{3.  **Numerical Comparison**:} \\
\texttt{    *   \$P(Kwoz | Kwox) = 55\textbackslash{}\%\$} \\
\texttt{    *   \$P(Kwoz | \textbackslash{}neg Kwox) = 34\textbackslash{}\%\$} \\
\\ 
\texttt{    Calculating the difference:} \\
\texttt{    \$\$55\textbackslash{}\% {-} 34\textbackslash{}\% = 21\textbackslash{}\%\$\$} \\
\\ 
\texttt{4.  **Conclusion**:} \\
\texttt{    Since the probability of \$Kwoz\$ occurring when \$Kwox\$ is present (55\%) is strictly greater than the probability when \$Kwox\$ is absent (34\%), the presence of \$Kwox\$ causally increases the likelihood of \$Kwoz\$. The direct pathways from \$Kwox\$ to \$Muvy\$ and \$Swoq\$ (which both lead to \$Kwoz\$) transmit this positive effect.} \\
\\ 
\texttt{Therefore, \$Kwox\$ increases the chance of \$Kwoz\$.} \\
            }
        }
    }
    & & \\ \\

    \theutterance \stepcounter{utterance}  
    & & & \multicolumn{2}{p{0.3\linewidth}}{
        \cellcolor[rgb]{0.9,0.9,0.9}{
            \makecell[{{p{\linewidth}}}]{
                \texttt{\tiny{[GM$|$GM]}}
                \texttt{Yes} \\
            }
        }
    }
    & & \\ \\

    \theutterance \stepcounter{utterance}  
    & & & \multicolumn{2}{p{0.3\linewidth}}{
        \cellcolor[rgb]{0.9,0.9,0.9}{
            \makecell[{{p{\linewidth}}}]{
                \texttt{\tiny{[GM$|$GM]}}
                \texttt{yes} \\
            }
        }
    }
    & & \\ \\

    \theutterance \stepcounter{utterance}  
    & & & \multicolumn{2}{p{0.3\linewidth}}{
        \cellcolor[rgb]{0.9,0.9,0.9}{
            \makecell[{{p{\linewidth}}}]{
                \texttt{\tiny{[GM$|$GM]}}
                \texttt{game\_result = WIN} \\
            }
        }
    }
    & & \\ \\

\end{supertabular}
}

\end{document}
